\reportheading{6. Lighting controls and energy use}

\subhead{6.1 Local and automatic control}

The office is divided into 32 independent control zones: 28 for bench groups and four for circulation, each with local manual switching and dimming. The largest zone measures 46.67 m², which is below the 250 m² automatic-control limit. According to BEC Table 5.5, the 1042.10 m² office requires 32 manual control points without needing a reduction for low power density [8], and Figure 5 illustrates their boundaries.

Each zone switches off after 10 minutes without occupancy, which falls within the BEC limit of 15 minutes. In addition, a schedule switches lighting off outside working hours, with local overrides limited to 120 minutes. To verify compliance, commissioning must confirm at least 50\% power reduction, including sensor and driver standby consumption [8, Sec. 5.6.1].

\reportfigure{\includegraphics[width=\linewidth]{zones.pdf}}{Figure 5. Proposed independent lighting control zones. The DIALux calculation uses one luminaire group.}

\subhead{6.2 Daylight response}

In areas where the external glazing and lighting load meet BEC clause 5.6.2 criteria, daylight-responsive control is proposed [8]. However, window dimensions and grouping are still needed to define these zones and check the exceptions. Where it is required, the controller must also be able to reduce lighting power to 50\% or less, or switch it off entirely, in response to daylight.

\subhead{6.3 Connected power}

The fittings draw 186 × 23 = 4278 W. With a 50 W allowance for sensors and shared control equipment, the design load reaches 4328 W, or 4.153 W/m² across 1042.10 m², which is below the BEC limit of 7.2 W/m². At full load for 2500 h/year, the energy use is about 10.82 MWh/year. However, a controlled estimate requires logged occupancy, driver power at dimmed output, and standby loads.
